% Theme: Poster  --  best for one concept per slide, quick definitions
% Build:  ./build.sh example-poster.tex   (or compile on Overleaf)
\documentclass[12pt]{beamer}
\usetheme{igposter}

\ighandle{@quantum_danesh}
\igseries{Qubits desde cero}
\igissue{03}

\begin{document}

\coverslide{Conceptos en 30 segundos}
  {CPTP}
  {Las cuatro letras detrás de todo canal cuántico.}

\begin{frame}{Respeta las mezclas}{Lineal}
{\LARGE
\begin{align*}
  &\Phi\big(p\rho + (1-p)\sigma\big)\\
  &\quad= p\,\Phi(\rho) + (1-p)\,\Phi(\sigma)
\end{align*}}%
{\large Mezclar y luego evolucionar es lo mismo que evolucionar y luego mezclar.\par}
\end{frame}

\begin{frame}{Estados a estados}{Positivo}
{\huge
\[
  \rho \succeq 0 \;\Longrightarrow\; \Phi(\rho) \succeq 0
\]}%
{\large Una matriz densidad nunca puede acabar dando \hl{probabilidades negativas}.\par}
\end{frame}

\begin{frame}{También con compañía}{Completamente positivo}
{\huge
\[
  (\mathrm{id}_k \otimes \Phi)(\rho) \succeq 0 \quad \forall k
\]}%
{\large La positividad debe sobrevivir cuando $\Phi$ actúa sobre \hl{la mitad
de un estado entrelazado}.\par}
\begin{notebox}
La transpuesta es positiva pero no completamente positiva. Por eso no es un
canal.
\end{notebox}
\end{frame}

\begin{frame}{La probabilidad total es 1}{Preserva la traza}
{\huge
\[
  \Tr \Phi(\rho) = \Tr \rho
\]}%
{\large Lo que entra normalizado sale normalizado.\par}
\end{frame}

\begin{frame}{La forma de Kraus}{Todo junto}
\keyeq{\Phi(\rho) = \sum_i K_i\,\rho\,K_i^\dagger,
  \qquad \sum_i K_i^\dagger K_i = \id}
Un mapa es CPTP \key{si y solo si} se puede escribir así.
\begin{defbox}[Teorema de Choi y Kraus]
Los $K_i$ no son únicos: dos conjuntos relacionados por una matriz unitaria
dan el mismo canal.
\end{defbox}
\end{frame}

\closingslide{Cuatro letras, un canal}
  {Mañana: por qué la transpuesta suspende el examen.}

\end{document}
